% Source: https://artofproblemsolving.com/community/c3h3166012_2019_lmt_spring_guts_round_p13p24__lexington_mathematical_tournament
% Stable identity: aops-topic-3166012-q22
The homothety centered at $B$ with ratio $-2$ sends $A$ to $D$ and $C$ to $E$, so it sends the first circumcircle of radius $R$ to the second of radius $2R$. Since both circles pass through $B$, their centers are on opposite sides of $B$ and are $3R$ apart. The external common-tangent segment therefore has length $\sqrt{(3R)^2-(2R-R)^2}=2\sqrt2R$. Heron's formula gives the original triangle area $\dfrac{15\sqrt7}{4}$, so $R=\dfrac{4\cdot5\cdot6}{15\sqrt7}=\dfrac{8}{\sqrt7}$. Hence $\boxed{MN=\dfrac{16\sqrt{14}}{7}}$.
\begin{asy}
usepackage("amsmath");
size(7cm);
pair Bp=(0,0),O1=(-1,0),O2=(2,0),u=(-1/3,sqrt(8)/3),M=O1+u,Np=O2+2*u;draw(circle(O1,1));draw(circle(O2,2));draw(M--Np,blue);draw(O1--M,dashed);draw(O2--Np,dashed);draw(O1--O2,gray);label("$B$",Bp,S);label("$O_1$",O1,S);label("$O_2$",O2,S);label("$M$",M,N);label("$N$",Np,N);
\end{asy}
